\chapter{Introducción}
\label{cap:introduccion}

\chapterquote{Yoshiii}{Rocío}

En los últimos años, se han presentado casos de personas que afirmaban poder jugar a un juego con tan solo un casco y su mente, sin pulsar botones. No obstante, estos artículos proporcionan escasa información al respecto, resultando poco fiables. En caso de que jugar a un juego ``con la mente'' sea posible, supondría un gran avance en los sistemas de accesibilidad, además de un enfoque novedoso y atractivo para el \textit{gaming}. 

A pesar de la existencia de dichos ejemplos, las interfaces cerebro-computador (\textit{Brain-Computer Interfaces}, BCI) han experimentado un crecimiento significativo real. Se ha avanzado tanto en investigación, consiguiendo modelos de inteligencia artificial (IA) con resultados muy prometedores, como en aplicaciones relacionadas con la accesibilidad y la interacción humano-computador. El progreso visible es alentador, sin embargo existe un problema: todos estos avances suelen necesitar de un gran presupuesto o cascos de lectura EEG poco manejables.

El presente trabajo de fin de grado se centra en el diseño e implementación de un sistema capaz de interpretar señales EEG en tiempo real y traducirlas en acciones utilizables dentro de distintos videojuegos. Para ello, se plantea desarrollar una arquitectura modular orientada a la adquisición continua de señales, al procesamiento de las mismas, la clasificación de la entrada y la traducción de la clasificación a una acción en tiempo real en un videojuego.

La propuesta plantea distintos retos debido a la naturaleza de las señales EEG: la presencia de ruido en las ondas, la variabilidad de las mismas y la dificultad de obtener predicciones estables en tiempo real, entre otros.

Además del propio \textit{pipeline} de procesamiento, se plantea diseñar diferentes juegos específicamente adaptados a las limitaciones y características del sistema BCI propuesto. Estos entornos permitirán evaluar el comportamiento del sistema en situaciones dinámicas e interactivas, facilitando tanto el análisis del rendimiento técnico como la experiencia de uso del usuario.

\section{Motivación}
Introducción al tema del TFG.


\section{Objetivos}
Descripción de los objetivos del trabajo.


\section{Plan de trabajo}
Aquí se describe el plan de trabajo a seguir para la consecución de los objetivos descritos en el apartado anterior.



\section{Explicaciones adicionales sobre el uso de esta plantilla}
Si quieres cambiar el \textbf{estilo del título} de los capítulos del documento, edita el fichero \verb|TeXiS\TeXiS_pream.tex| y comenta la línea \verb|\usepackage[Lenny]{fncychap}| para dejar el estilo básico de \LaTeX.

Si no te gusta que no haya \textbf{espacios entre párrafos} y quieres dejar un pequeño espacio en blanco, no metas saltos de línea (\verb|\\|) al final de los párrafos. En su lugar, busca el comando  \verb|\setlength{\parskip}{0.2ex}| en \verb|TeXiS\TeXiS_pream.tex| y aumenta el valor de $0.2ex$ a, por ejemplo, $1ex$.

TFGTeXiS se ha elaborado a partir de la plantilla de TeXiS\footnote{\url{http://gaia.fdi.ucm.es/research/texis/}}, creada por Marco Antonio y Pedro Pablo Gómez Martín para escribir su tesis doctoral. Para explicaciones más extensas y detalladas sobre cómo usar esta plantilla, recomendamos la lectura del documento \texttt{TeXiS-Manual-1.0.pdf} que acompaña a esta plantilla.

El siguiente texto se genera con el comando \verb|\lipsum[2-20]| que viene a continuación en el fichero .tex. El único propósito es mostrar el aspecto de las páginas usando esta plantilla. Quita este comando y, si quieres, comenta o elimina el paquete \textit{lipsum} al final de \verb|TeXiS\TeXiS_pream.tex|

\subsection{Texto de prueba}


\lipsum[2-20]